% Generated draft: requires mathematical and source review.
\documentclass[11pt]{article}
\usepackage[margin=25mm]{geometry}
\usepackage{fontspec}
\setmainfont{Latin Modern Roman}
\usepackage{amsmath,amssymb,mathtools,unicode-math}
\setmathfont{Latin Modern Math}
\usepackage{enumitem,needspace,xurl,hyperref}
\setlength{\parindent}{0pt}
\setlength{\parskip}{5pt}
\setlength{\emergencystretch}{3em}
\newcommand{\N}{\mathbb{N}}
\newcommand{\Z}{\mathbb{Z}}
\newcommand{\Q}{\mathbb{Q}}
\newcommand{\R}{\mathbb{R}}
\newcommand{\C}{\mathbb{C}}
\newcommand{\F}{\mathbb{F}}
\newcommand{\A}{\mathbb{A}}
\newcommand{\PP}{\mathbb{P}}
\newcommand{\E}{\mathbb{E}}
\newcommand{\eps}{\varepsilon}
\newcommand{\dd}{\,\mathrm{d}}
\begin{document}
\section*{UnsolvedMath ID: 701577}
\textbf{Source problem code:} catalogue entry
\subsection{Definitions and mathematical statement}
Let \(G\) be a group. Write \(\operatorname{Aut}(G)\) for its automorphism group and \(\operatorname{Inn}(G)=\{\widehat g:g\in G\}\), where \(\widehat g(x)=gxg^{-1}\). Set
\[
\operatorname{Out}(G)=\operatorname{Aut}(G)/\operatorname{Inn}(G).
\]
Regard an element \(\Phi\in\operatorname{Out}(G)\) as its coset of representatives in \(\operatorname{Aut}(G)\). Two representatives \(\varphi,\psi\in\Phi\) are called \emph{isogradient} if
\[
\varphi=\widehat g^{-1}\circ\psi\circ\widehat g
\qquad\text{for some }g\in G.
\]
Let \(S(\Phi)\) be the cardinality of the set of equivalence classes under this relation.

A group is \emph{residually finite} if, for each \(x\in G\setminus\{1\}\), there exist a finite group \(F\) and a homomorphism \(q:G\to F\) such that \(q(x)\ne1\). A group \(H\) is \emph{soluble} if its derived series, defined by \(H^{(0)}=H\) and \(H^{(n+1)}=[H^{(n)},H^{(n)}]\), satisfies \(H^{(d)}=\{1\}\) for some integer \(d\ge0\).

\textbf{Conjecture (A. L. Fel'shtyn and E. V. Troitsky; Kourovka Notebook, Problem 19.29).} For every finitely generated residually finite group \(G\),
\[
\bigl(\exists\,\Phi\in\operatorname{Out}(G)\text{ with }S(\Phi)<\infty\bigr)
\ \Longrightarrow\ 
\bigl(\exists\,H\le G\text{ with }[G:H]<\infty
\text{ and }H\text{ soluble}\bigr).
\]
The supplied registry labels this conjecture open; that label is supplied source information and is not an independently verified assessment.
\subsection{Short English statement}
The conjecture asks whether a finitely generated residually finite group must contain a soluble subgroup of finite index whenever some outer automorphism class contains only finitely many classes of representatives under conjugation by inner automorphisms.
\subsection{Sources}
\begin{itemize}
\item Supplied local UnsolvedMath registry record: permanent ID \(701577\), exact source code \texttt{catalogue entry}, title \emph{Kourovka Notebook Problem 19.29}; upstream revision \texttt{d371ee917ac9647fbb06cd6fc529deec6f3dda82}.
\item \emph{The Kourovka Notebook (Unsolved Problems in Group Theory)}, Problem 19.29, attributed in the supplied record to A. L. Fel'shtyn and E. V. Troitsky, proposed in 2018. The supplied background identifies the 19th issue (2018) and the 21st issue edition (2026), \texttt{arXiv:1401.0300v48}. The collection description names E. I. Khukhro and V. D. Mazurov as editors.
\end{itemize}
\end{document}
